\documentclass{beamer}
\usepackage{ctex}
\usetheme{SimplePlus}
\setbeamertemplate{footline}{}

\usepackage{amsmath,amssymb,mathtools}
\usepackage{tikz}
\usetikzlibrary{arrows.meta}
\usepackage{tikz-cd}

\usepackage{unicode-math}
%\setmathfont{STIX Two Math}
%\setmathfont{TeX Gyre Termes Math}
%\setmathfont{Libertinus Math}
%\setmathfont{Fira Math}

\usepackage[
  hybrid,                            % 允许在 Markdown 中混用 LaTeX 命令
  texMathDollars                     % 启用 $...$ 和 $$...$$ 数学公式
]{markdown}

\usepackage{graphicx} % Allows including images
\graphicspath{{./images/}}

\title{presheaf}

\date{}

\begin{document}

\begin{frame}
    % Print the title page as the first slide
    \titlepage
\end{frame}

\begin{frame}[fragile]{子函子(Subfunctor)的定义}
\begin{markdown}
设
$$
F,G:\mathcal C\to\mathcal D
$$

是两个函子。

如果对每个 $X\in\mathcal C$，都有一个 monomorphism
$$
i_X:G(X)\hookrightarrow F(X),
$$

并且这些 $i_X$ 组成一个自然变换
$$
i:G\Rightarrow F,
$$

那么称
$$
\boxed{G\text{ 是 }F\text{ 的一个 subfunctor}}
$$

或者说
$$
G\subseteq F.
$$

换句话说：
$$
\boxed{\text{subfunctor}=\text{objectwise subobject}+\text{naturality}}
$$
\end{markdown}
\end{frame}

\begin{frame}[fragile]{筛(Sieve)}
\begin{markdown}
给定 $X\in\mathcal C$。

一个 sieve $S$ on $X$ 是一个 subfunctor
$$ S\subseteq y_X=\mathcal C(-,X). $$

等价地：

$S$ 是一类指向 $X$ 的 morphisms，并且对于复合封闭。

即：

如果
$$ f:Y\to X\in S $$

以及
$$ g:Z\to Y, $$

那么
$$ f\circ g:Z\to X $$

也必须属于 $S$。
\end{markdown}
\end{frame}

\begin{frame}[fragile]{Grothendieck topology 的定义}
\begin{markdown}
一个 Grothendieck topology $J$ 给每一个对象
$$
X\in\mathcal C
$$

指定一个集合
$$
J(X)
$$

其中元素都是 $X$ 上的 sieves：
$$
J(X)\subseteq\operatorname{Sieve}(X).
$$

如果
$$
S\in J(X),
$$

我们就称：
$$
\boxed{S\text{ 是 }X\text{ 上的 covering sieve}}
$$

所以：
$$
\boxed{J(X)=\\{\text{$X$ 上所有被认为是 covering 的 sieves}\\}}
$$
\end{markdown}
\end{frame}

\begin{frame}[fragile]
\begin{markdown}
一个 **Grothendieck topology** $J$ on $\mathcal C$，是一个赋值
$$
X\mapsto J(X)
$$

满足下面三个公理。

---

\textbf{公理 GT1：最大 sieve 是 covering}

对每个 $X$：
$$
\boxed{
y_X\in J(X)
}
$$

这里
$$
y_X=\mathcal C(-,X)
$$

是最大 sieve。

为什么它叫最大 sieve？

因为对于任何 sieve $S$ on $X$：
$$
S\subseteq y_X.
$$

所以 $y_X$ 包含所有
$$
Y\to X.
$$
\end{markdown}
\end{frame}

\begin{frame}[fragile]
\begin{markdown}
\textbf{公理 GT2：Base change / Pullback stability}

假设：
$$
S\in J(X)
$$

即 $S$ 是 $X$ 上的 covering sieve。

再取任意 morphism：
$$
f:Y\to X.
$$

我们可以把 $S$ 拉回到 $Y$。

定义：
$$
\boxed{f^*S = \{g:Z\to Y\mid f\circ g\in S\}}
$$

那么要求：
$$
\boxed{S\in J(X)\quad\Longrightarrow\quad f^*S\in J(Y)}
$$

也就是说：

> **cover 拉回以后仍然是 cover。**
\end{markdown}
\end{frame}

\begin{frame}[fragile]
\begin{markdown}
\textbf{公理 GT3：Transitivity}

这是三个公理中最不直观、但非常重要的一个。

假设：
$$
S\in J(X).
$$

所以 $S$ 是 $X$ 的一个 covering sieve。

现在对于每个
$$
f:Y\to X
$$

且
$$
f\in S,
$$

我们又选择一个 covering sieve：
$$
R_f\in J(Y).
$$

把所有这些 $R_f$ “合起来”，得到 $X$ 上的 sieve $R$：
$$
R=
\{g:Z\to X
\mid
\exists f:Y\to X\in S,\;
\exists h:Z\to Y\in R_f,\;
g=f\circ h
\}.
$$

那么要求：
$$
\boxed{
R\in J(X)
}
$$

这就是 transitivity。
\end{markdown}
\end{frame}

\begin{frame}[fragile]
\begin{markdown}
\textbf{用一句话理解 GT3}

它说：
$$
\boxed{
\text{cover }X +
\text{cover each piece}
\Rightarrow
\text{cover }X
}
$$

例如：
$$
X
$$

先被
$$
U_i\to X
$$

覆盖。

然后每一个 $U_i$ 又被
$$
V_{ij}\to U_i
$$

覆盖。

那么复合：
$$
V_{ij}\to U_i\to X
$$

应该仍然构成 $X$ 的覆盖。

这就是所谓：
$$
\boxed{\text{cover of a cover is a cover}}
$$
\end{markdown}
\end{frame}

\begin{frame}[fragile]{Grothendieck topology 另一个版本定义(貌似更常用)}
\begin{markdown}
- 一个 **Grothendieck 拓扑** $J$（或称 coverage / pretopology 的饱和化）在 $\mathcal{C}$ 上给每个对象 $X$ 指定一族 **covering sieves**（覆盖筛）$J(X)$，满足：  
  1. 最大筛 $y(X)=\mathrm{Hom}(-,X)$ 属于 $J(X)$；  
  2. 对任意 $f:Y\to X$，若 $R\in J(X)$，则拉回 $f^*R\in J(Y)$；  
  3. 若 $R\in J(X)$，且筛 $S$ 满足“对每个 $f\in R$ 有 $f^*S\in J(\mathrm{dom}\,f)$”，则 $S\in J(X)$。  
- **位点** $(\mathcal{C},J)$ 就是配备了 Grothendieck 拓扑的范畴。  
（常用的“覆盖族”版本要求 pullback 存在，并满足同构、复合、拉回三条公理；再取生成的筛即得拓扑。）
\end{markdown}
\end{frame}

\begin{frame}[fragile]{Site}
\begin{markdown}
$$
\boxed{
\text{Site} =
(\mathcal C,J)
}
$$

其中：

* $\mathcal C$：一个范畴；
* $J$：$\mathcal C$ 上的 Grothendieck topology。

这就是 site。

注意：

> **site 不是一个新的数学对象类别，而是“一个范畴 + 指定哪些 sieve 是覆盖”的结构。**
\end{markdown}
\end{frame}

\begin{frame}[fragile]{$S$-compatible family}
\begin{markdown}
我们有：

$$
S\hookrightarrow y_X
$$

和 presheaf
$$
F:\mathcal C^{op}\to\mathbf{Set}.
$$

一个 $S$-compatible family 定义为一个自然变换：
$$
\boxed{
s:S\Rightarrow F.
}
$$

也就是说：
$$
s\in\operatorname{Nat}(S,F).
$$
\end{markdown}
\end{frame}

\begin{frame}[fragile,shrink]
\begin{markdown}
因为：
$$
S(Y)\subseteq\mathcal C(Y,X),
$$

所以对于每个：
$$
f:Y\to X
$$

且
$$
f\in S(Y),
$$

我们得到一个：
$$
s_f\in F(Y).
$$

也就是说：
$$
\boxed{
\text{每一个 cover 的局部箭头 }f:Y\to X
\text{ 都有一个 section }s_f\in F(Y)
}
$$

但是这些 $s_f$ 不能随便选。

自然性要求：

如果
$$
g:Z\to Y,
$$

那么：
$$
f\circ g\in S(Z),
$$

并且要求：
$$
\boxed{
F(g)(s_f)=s_{f\circ g}.
}
$$

也就是说：

> 把 $s_f$ 限制到更小的 $Z$，必须得到 $s_{f\circ g}$。

所以它确实是一个**相容的局部数据系统**。
\end{markdown}
\end{frame}

\begin{frame}[fragile,shrink]{一个 global section 如何产生 compatible family？}
\begin{markdown}
现在假设我们已经有：

$$
s\in F(X).
$$

那么对每一个：

$$
f:Y\to X\in S,
$$

可以把 $s$ restriction 到 $Y$：

$$
F(f)(s)\in F(Y).
$$

因此：

$$
s\in F(X)
$$

自然地产生一个自然变换：

$$
S\Rightarrow F.
$$

于是有一个 canonical map：

$$
\boxed{
F(X)\longrightarrow\operatorname{Nat}(S,F).
}
$$

它可以理解成：

$$
\boxed{
\text{global section}
\longmapsto
\text{its compatible local restrictions}.
}
$$
\end{markdown}
\end{frame}

\begin{frame}[fragile]{Sheaf condition}
\begin{markdown}
一个 presheaf
$$
F:\mathcal C^{op}\to\mathbf{Set}
$$

称为一个 **sheaf**，如果对于每一个对象
$$
X\in\mathcal C
$$

以及每一个 covering sieve
$$
S\in J(X),
$$

canonical map
$$
\boxed{
F(X)\longrightarrow\operatorname{Nat}(S,F)
}
$$

都是双射。

即：
$$
\boxed{
F(X)\cong\operatorname{Nat}(S,F)
}
$$

这就是 **sieve formulation of the sheaf condition**。
\end{markdown}
\end{frame}

\begin{frame}[fragile]{这个双射究竟表达了什么？}
\begin{markdown}
它实际上包含两个条件。

---

\textbf{(1) Existence：可以 glue}

surjective：

$$
F(X)\twoheadrightarrow\operatorname{Nat}(S,F)
$$

意味着：

> 每一个 compatible local family 都来自某个 global section。

即：

$$
\boxed{
\text{local compatible data}
\Rightarrow
\text{global data}
}
$$

这就是 **gluing**。
\end{markdown}
\end{frame}

\begin{frame}[fragile]
\begin{markdown}
\textbf{Uniqueness：gluing 是唯一的}

injective：
$$
F(X)\hookrightarrow\operatorname{Nat}(S,F)
$$

意味着：

如果
$$
s,t\in F(X)
$$

在 cover $S$ 上 restriction 完全相同，那么：
$$
s=t.
$$

即：
$$
\boxed{
\text{global section is determined locally}
}
$$

所以 sheaf condition 实际上就是：
$$
\boxed{
\text{existence of gluing}
+
\text{uniqueness of gluing}
}
$$
\end{markdown}
\end{frame}

\begin{frame}[fragile]
\begin{markdown}
从 Yoneda 观点看：
$$
S\hookrightarrow y_X.
$$

restriction 给出：
$$
F(X)
\cong
\operatorname{Nat}(y_X,F)
\longrightarrow
\operatorname{Nat}(S,F).
$$

因此 sheaf condition 就是：
$$
\boxed{
\operatorname{Nat}(y_X,F)
\simeq
\operatorname{Nat}(S,F)
}
$$

这已经开始显露出 **localization** 的味道：

> 对于 sheaf $F$，cover $S\to y_X$ 在 $F$ 看来是一个等价。
\end{markdown}
\end{frame}

\begin{frame}[fragile]{Presheaf category 与 sheaf category}
\begin{markdown}
所有 presheaves：
$$
\widehat{\mathcal C} = [\mathcal C^{op},\mathbf{Set}].
$$

所有 sheaves：
$$
\operatorname{Sh}(\mathcal C,J) \subseteq \widehat{\mathcal C}.
$$

这是一个非常特殊的子范畴：
$$
\boxed{\operatorname{Sh}(\mathcal C,J) \hookrightarrow \operatorname{PSh}(\mathcal C)}
$$

具有左伴随：
$$
a:\operatorname{PSh}(\mathcal C)\rightleftarrows\operatorname{Sh}(\mathcal C,J):i
$$

其中 $i$ 是 inclusion，$a$ 称为 **sheafification**。

所以：
$$
a\dashv i.
$$
\end{markdown}
\end{frame}

\begin{frame}[fragile]{Sheafification}
\begin{markdown}
对于任意 presheaf $F$，sheafification 给出
$$
aF
$$

它是一个 sheaf，并且有 canonical map
$$
\eta_F:F\to i(aF)
$$
\end{markdown}
\end{frame}

\begin{frame}[fragile,shrink]{filtered category}
\begin{markdown}
**Filtered category（滤子化范畴 / 滤范畴）的严格定义如下：**

一个范畴 $J$ 称为 **filtered（滤子化的）**，如果它满足以下三个条件：

1. $J$ 非空；
2. 对任意两个对象 $j, j' \in J$，存在对象 $k \in J$ 以及态射 $j \to k$ 和 $j' \to k$；
3. 对任意两个平行态射 $u, v : i \to j$，存在对象 $k$ 以及态射 $w : j \to k$，使得 $w \circ u = w \circ v$。

这是最常见的“有限滤子化”（finitely filtered）定义，也是 Wikipedia、Stacks Project、nLab 等标准来源采用的版本。

等价的刻画（更抽象、更统一）是：

- $J$ 中**任意有限图（finite diagram）都有 cocone**（即存在某个对象以及从该图到这个对象的自然变换）。

这两者等价，因为任意有限 colimit 可由初始对象、二元余积和余等化子生成，而上述三条条件正好保证这三类有限图有 cocone。

（更一般地，对正则基数 $\kappa$，可定义 $\kappa$-filtered category：任意基数小于 $\kappa$ 的图都有 cocone。普通 filtered 对应 $\omega$-filtered / 有限情况。）

\textbf{直观解释}

Filtered category 是 **directed set（有向集）** 的范畴化版本。

- 在有向集中，任意两个元素都有一个“共同上界”。
- 在范畴中，对象之间没有“大小比较”，只有态射，因此：
  - 条件 2 说：任意两个对象都能“往后走”找到一个共同的目标对象（上界）；
  - 条件 3 说：如果两个平行态射“不一致”，你还可以再往后走一步，把它们“强行变成相等”（等于是用一个态射把它们等化掉）。

整体感觉是：这个范畴“越往后走越集中、越一致”，你可以不断往后推进，把任意有限多的对象和态射都“汇聚”到同一个对象上，并且让相关的三角形交换。

**为什么重要？**  
以 filtered category 为指标的 colimit（称为 **filtered colimit**）行为特别好：它们在 $\mathbf{Set}$ 中与有限极限交换，很多构造（如 Ind-对象、局部有限表示模等）都依赖这个性质。

对偶概念是 **cofiltered category**（上滤子化范畴）：把所有箭头方向反过来即可。
\end{markdown}
\end{frame}

\begin{frame}[fragile,shrink]{pseudofunctor}
现在我们考虑：
$$
\mathcal F:\mathcal C^{op}\to\mathbf{Cat}.
$$

对于每个：
$$
X\in\mathcal C
$$

给一个 category：
$$
\mathcal F(X).
$$

对于每个：
$$
f:Y\to X
$$

给一个 functor：
$$
f^*:\mathcal F(X)\to\mathcal F(Y).
$$

我们希望：
$$
(f\circ g)^*=g^*\circ f^*
$$

并且：
$$
(\operatorname{id}_X)^*=\operatorname{id}_{\mathcal F(X)}.
$$

但是在很多自然例子中，这两个等式并不是严格成立，而只有\textbf{自然同构}：
$$
\boxed{(f\circ g)^*\cong g^*\circ f^*}
$$

以及：
$$
\boxed{(\operatorname{id}_X)^*\cong\operatorname{id}_{\mathcal F(X)}}
$$

因此我们不能用普通 functor，而使用 \textbf{pseudofunctor}。
\end{frame}

\begin{frame}[fragile,shrink]{Composition constraint}
如果：
$$
X\xleftarrow{f}Y\xleftarrow{g}Z
$$

即：
$$
f:Y\to X
\qquad
g:Z\to Y
$$

那么有：
$$
g^*:\mathcal F(Y)\to\mathcal F(Z)
$$

和：
$$
f^*:\mathcal F(X)\to\mathcal F(Y).
$$

因此：
$$
g^*f^*:\mathcal F(X)\to\mathcal F(Z).
$$

另一方面：
$$
f\circ g:Z\to X
$$

给出：
$$
(f\circ g)^*:\mathcal F(X)\to\mathcal F(Z).
$$

pseudofunctor 给出一个自然同构：
$$
\boxed{
c_{f,g}:
g^*f^*
\xRightarrow{\cong}
(f\circ g)^*
}
$$

注意这里的 $c_{f,g}$ 是两个 functor 之间的 \textbf{natural isomorphism}。
\end{frame}

\begin{frame}[fragile]{Identity constraint}
对于每个：
$$
X\in\mathcal C,
$$

有：
$$
\operatorname{id}_X:X\to X.
$$

因此有：
$$
(\operatorname{id}_X)^*:
\mathcal F(X)\to\mathcal F(X).
$$

pseudofunctor 给出自然同构：
$$
\boxed{
u_X:
\operatorname{Id}_{\mathcal F(X)}
\xRightarrow{\cong}
(\operatorname{id}_X)^*.
}
$$

当然也可以把方向反过来定义：
$$
(\operatorname{id}_X)^*
\xRightarrow{\cong}
\operatorname{Id}_{\mathcal F(X)}.
$$

两种 convention 都可以，关键是后面的 coherence condition 必须与 convention 一致。
\end{frame}

\begin{frame}[fragile]{Pseudofunctor 还必须满足 coherence}
\begin{markdown}
仅仅给：

$$
c_{f,g}
$$

和：

$$
u_X
$$

还不够。

它们必须满足两个 coherence conditions：

1. identity coherence；
2. associativity coherence。

这就是 pseudofunctor 的核心。
\end{markdown}
\end{frame}

\begin{frame}[fragile,shrink]{Associativity coherence}
考虑三个 composable arrows：
$$
W\xrightarrow{h}Z\xrightarrow{g}Y\xrightarrow{f}X
$$

即：
$$
f:Y\to X
\qquad
g:Z\to Y
\qquad
h:W\to Z
$$

我们有四个相关 functors：
$$
h^*g^*f^*,
$$

$$
(h\circ g)^*f^*,
$$

$$
h^*(f\circ g)^*,
$$

以及：
$$
(f\circ g\circ h)^*.
$$

从：
$$
h^*g^*f^*
$$

到：
$$
(fgh)^*
$$

存在两条自然同构路径。
\end{frame}

\begin{frame}[fragile]
\textbf{路径 A：先组合 $f,g$}

首先：
$$
h^*g^*f^*
\xRightarrow{h^**c_{f,g}}
h^*(fg)^*.
$$

然后：
$$
h^*(fg)^*
\xRightarrow{c_{fg,h}}
(fgh)^*.
$$

所以得到：
$$
h^*g^*f^*
\Longrightarrow
(fgh)^*.
$$
\end{frame}

\begin{frame}[fragile]
\begin{markdown}
**路径 B：先组合 $g,h$**

首先：
$$
h^*g^*f^*
\xRightarrow{c_{g,h}*f^*}
(gh)^*f^*.
$$

然后：
$$
(gh)^*f^*
\xRightarrow{c_{f,gh}}
(fgh)^*.
$$

要求这两条路径相等：
$$
\boxed{
c_{fg,h}\circ(h^**c_{f,g}) =
c_{f,gh}\circ(c_{g,h}*f^*).
}
$$

这里：

* $h^**c_{f,g}$ 表示把自然变换 $c_{f,g}$ 用 functor $h^*$ whisker；
* $c_{g,h}*f^*$ 表示在右边 whisker $f^*$。

这就是 pseudofunctor 的 **associativity coherence condition**。
\end{markdown}
\end{frame}

\begin{frame}[fragile,shrink]{whiskering}
先固定 whiskering 记号

这是前面最容易混淆的地方。

假设：
$$
\alpha:F\Rightarrow G
$$

是 natural transformation。

\textbf{左 whiskering}

如果：
$$
H:\mathcal B\to\mathcal D,
$$

那么：
$$
\boxed{
H*\alpha
}
$$

定义为：
$$
HF\Rightarrow HG.
$$

其 component：
$$
\boxed{
(H*\alpha)_A
=
H(\alpha_A).
}
$$

\textbf{右 whiskering}

如果：
$$
K:\mathcal D\to\mathcal A,
$$

那么：
$$
\boxed{
\alpha*K
}
$$

定义为：

$$
FK\Rightarrow GK.
$$

其 component：
$$
\boxed{
(\alpha*K)_D
=
\alpha_{K(D)}.
}
$$

\textbf{Vertical composition}

如果：
$$
F\xRightarrow{\alpha}G
\xRightarrow{\beta}H,
$$

则：
$$
\boxed{
\beta\circ\alpha:F\Rightarrow H
}
$$

是 vertical composition。

其 component：
$$
\boxed{
(\beta\circ\alpha)_A
=
\beta_A\circ\alpha_A.
}
$$

因此：
$$
\circ
$$

和：
$$
*
$$

是两种不同的操作：
$$
\boxed{
\circ=\text{vertical composition},
\qquad
*=\text{whiskering}.
}
$$
\end{frame}

\begin{frame}[fragile,shrink]{Right identity coherence}
考虑：
$$
f:Y\to X.
$$

composition constraint：
$$
c_{f,\operatorname{id}_Y}
:
(\operatorname{id}_Y)^*f^*
\Longrightarrow
f^*.
$$

另一方面：
$$
u_Y^{-1}:
(\operatorname{id}_Y)^*
\Longrightarrow
\operatorname{Id}_{F(Y)}.
$$

右 whiskering \(f^*\)：
$$
u_Y^{-1}*f^*
:
(\operatorname{id}_Y)^*f^*
\Longrightarrow
f^*.
$$

要求：
$$
\boxed{
c_{f,\operatorname{id}_Y}
=
u_Y^{-1}*f^*.
}
$$

逐对象写就是：
$$
\boxed{
c_{f,\operatorname{id}_Y}(A)
=
u_Y^{-1}(f^*A).
}
$$
\end{frame}

\begin{frame}[fragile,shrink]{Left identity coherence}
同样：
$$
c_{\operatorname{id}_X,f}
:
f^*(\operatorname{id}_X)^*
\Longrightarrow
f^*.
$$

而：
$$
u_X^{-1}:
(\operatorname{id}_X)^*
\Longrightarrow
\operatorname{Id}_{F(X)}.
$$

左边 whiskering \(f^*\)：
$$
f^**u_X^{-1}
:
f^*(\operatorname{id}_X)^*
\Longrightarrow
f^*.
$$

所以要求：
$$
\boxed{
c_{\operatorname{id}_X,f}
=
f^**u_X^{-1}.
}
$$

逐对象：
$$
\boxed{
c_{\operatorname{id}_X,f}(A)
=
f^*\bigl(u_X^{-1}(A)\bigr).
}
$$
\end{frame}

\begin{frame}[fragile,shrink]{pseudofunctor定义总结}
一个 contravariant pseudofunctor
$$
\boxed{
F:\mathcal C^{op}\to\mathbf{Cat}
}
$$

包括：

\textbf{数据}
$$
\boxed{
\begin{aligned}
X&\mapsto F(X),\\
(f:Y\to X)&\mapsto f^*:F(X)\to F(Y),\\
c_{f,g}&:g^*f^*\xRightarrow{\cong}(fg)^*,\\
u_X&:\operatorname{Id}_{F(X)}
\xRightarrow{\cong}
(\operatorname{id}_X)^*.
\end{aligned}
}
$$

\textbf{条件}

对于每个 \(f:Y\to X\)：
$$
\boxed{
c_{f,\operatorname{id}_Y}
=
u_Y^{-1}*f^*
}
$$

以及：
$$
\boxed{
c_{\operatorname{id}_X,f}
=
f^**u_X^{-1}.
}
$$

对于：
$$
X\xleftarrow fY\xleftarrow gZ\xleftarrow hW,
$$

有：
$$
\boxed{
c_{fg,h}\circ(h^**c_{f,g})
=
c_{f,gh}\circ(c_{g,h}*f^*).
}
$$
\end{frame}

\begin{frame}[fragile]{用一个图看清 associativity}
可以把 associativity coherence 记成这个交换图：
$$
\boxed{
\begin{array}{ccc}
h^*g^*f^*
&
\xrightarrow{\,h^**c_{f,g}\,}
&
h^*(fg)^*
\\[2mm]
\downarrow{\scriptstyle c_{g,h}*f^*}
&&
\downarrow{\scriptstyle c_{fg,h}}
\\[2mm]
(gh)^*f^*
&
\xrightarrow{\,c_{f,gh}\,}
&
(fgh)^*
\end{array}
}
$$

即：
$$
\boxed{
c_{fg,h}\circ(h^**c_{f,g})
=
c_{f,gh}\circ(c_{g,h}*f^*).
}
$$
\end{frame}

\begin{frame}[fragile,shrink]{Descent datum}
设：
$$
\mathcal C
$$

是一个范畴，并且
$$
\boxed{
F:\mathcal C^{op}\to\mathbf{Cat}
}
$$

是一个 contravariant pseudofunctor。

我们固定以下记号。

对于：
$$
f:Y\to X,
$$

写：
$$
\boxed{
f^*:F(X)\to F(Y)
}
$$

对于：
$$
X\xleftarrow{f}Y\xleftarrow{g}Z,
$$

即：
$$
f:Y\to X,\qquad g:Z\to Y,
$$

pseudofunctor 的 composition constraint 是：
$$
\boxed{
c_{f,g}:
g^*f^*
\xRightarrow{\cong}
(f\circ g)^*
}
$$

对于 \(A\in F(X)\)，其 component 是：
$$
\boxed{
c_{f,g}(A):
g^*f^*A
\xrightarrow{\cong}
(fg)^*A
}
$$

Identity constraint 是：
$$
\boxed{
u_X:
\operatorname{Id}_{F(X)}
\xRightarrow{\cong}
(\operatorname{id}_X)^*
}
$$

所以：
$$
u_X(A):
A\xrightarrow{\cong}(\operatorname{id}_X)^*A
$$
\end{frame}

\begin{frame}[fragile,shrink]{固定一个 sieve}
现在取：
$$
X\in\mathcal C
$$

以及一个 sieve：
$$
\boxed{
S\hookrightarrow y_X=\mathcal C(-,X)
}
$$

也就是说，对于每一个 \(Y\in\mathcal C\)：
$$
S(Y)\subseteq\mathcal C(Y,X),
$$

并且：
$$
f\in S(Y),\quad g:Z\to Y
$$

必然推出：
$$
fg=f\circ g\in S(Z)
$$

如果：
$$
S\in J(X),
$$

则 \(S\) 是 covering sieve。

\textbf{Descent datum 的定义本身不需要 \(S\) 是 covering sieve；只有讨论 sheaf/stack condition 时才要求 \(S\) covering。}
\end{frame}

\begin{frame}[fragile]{Descent datum 的第一部分：局部对象}
对于每一个：
$$
f:Y\to X
$$

满足：
$$
f\in S(Y),
$$

选择一个对象：
$$
\boxed{
A_f\in F(Y)
}
$$

因此，我们为 sieve 中的每一个箭头：
$$
f:Y\to X
$$

配置一个位于 fiber \(F(Y)\) 中的对象 \(A_f\)。

可以写成：
$$
\boxed{
f\in S(Y)
\quad\leadsto\quad
A_f\in F(Y)
}
$$
\end{frame}

\begin{frame}[fragile,shrink]{Descent datum 的第二部分：transition isomorphism}
现在取：
$$
f:Y\to X\in S(Y)
$$

以及任意：
$$
g:Z\to Y
$$

因为 \(S\) 是 sieve，所以：
$$
fg=f\circ g\in S(Z)
$$

因此：
$$
A_f\in F(Y)
$$

可以沿着 \(g\) restriction：
$$
g^*A_f\in F(Z)
$$

另一方面：
$$
A_{fg}\in F(Z)
$$

因此要求给出一个 isomorphism：
$$
\boxed{
\theta_{f,g}:
g^*A_f
\xrightarrow{\cong}
A_{fg}
}
$$

这就是 descent datum 的 transition isomorphism。
\end{frame}

\begin{frame}[fragile,shrink]{第一个 coherence：identity condition}
现在令：
$$
g=\operatorname{id}_Y
$$

那么：
$$
fg=f
$$

所以：
$$
\theta_{f,\operatorname{id}_Y}:
(\operatorname{id}_Y)^*A_f
\longrightarrow
A_f
$$

另一方面，由 pseudofunctor 的 unit constraint：
$$
u_Y:
\operatorname{Id}_{F(Y)}
\Rightarrow
(\operatorname{id}_Y)^*
$$

有：
$$
u_Y(A_f):
A_f
\longrightarrow
(\operatorname{id}_Y)^*A_f
$$

因此要求：
$$
\boxed{
\theta_{f,\operatorname{id}_Y}
\circ
u_Y(A_f)
=
\operatorname{id}_{A_f}
}
$$

等价地：
$$
\boxed{
\theta_{f,\operatorname{id}_Y}
=
u_Y(A_f)^{-1}
}
$$

这就是 identity condition。
\end{frame}

\begin{frame}[fragile]{第二个 coherence：cocycle condition}
现在考虑三个 composable arrows：
$$
\boxed{
X\xleftarrow{f}Y\xleftarrow{g}Z\xleftarrow{h}W.
}
$$

即：
$$
f:Y\to X,
\qquad
g:Z\to Y,
\qquad
h:W\to Z.
$$

从：
$$
h^*g^*A_f
$$

到：
$$
A_{fgh}
$$

有两条自然的路径。
\end{frame}

\begin{frame}[fragile]{路径一：先沿 \(g\)，再沿 \(h\)}
首先使用：
$$
\theta_{f,g}:
g^*A_f\to A_{fg}.
$$

应用 \(h^*\)，得到：
$$
h^*(\theta_{f,g}):
h^*g^*A_f
\longrightarrow
h^*A_{fg}.
$$

然后使用：
$$
\theta_{fg,h}:
h^*A_{fg}
\longrightarrow
A_{fgh}.
$$

因此得到：
$$
\boxed{
h^*g^*A_f
\xrightarrow{h^*(\theta_{f,g})}
h^*A_{fg}
\xrightarrow{\theta_{fg,h}}
A_{fgh}.
}
$$
\end{frame}

\begin{frame}[fragile]{路径二：先把 \(g,h\) 合并}
pseudofunctor 给出：
$$
c_{g,h}:
h^*g^*
\xRightarrow{\cong}
(gh)^*.
$$

在 \(A_f\) 上：
$$
c_{g,h}(A_f):
h^*g^*A_f
\longrightarrow
(gh)^*A_f.
$$

然后使用：
$$
\theta_{f,gh}:
(gh)^*A_f
\longrightarrow
A_{fgh}.
$$

因此得到：
$$
\boxed{
h^*g^*A_f
\xrightarrow{c_{g,h}(A_f)}
(gh)^*A_f
\xrightarrow{\theta_{f,gh}}
A_{fgh}.
}
$$
\end{frame}

\begin{frame}[fragile]{Cocycle condition 的严格公式}
要求上面的两条路径相同：
$$
\boxed{
\theta_{fg,h}
\circ
h^*(\theta_{f,g})
=
\theta_{f,gh}
\circ
c_{g,h}(A_f).
}
$$

这就是 descent datum 的 cocycle condition。

注意这里的
$$
c_{g,h}(A_f)
$$

非常重要，因为 \(c_{g,h}\) 是一个 natural isomorphism：
$$
c_{g,h}:
h^*g^*
\Rightarrow
(gh)^*,
$$

而不是一个普通 morphism。
\end{frame}

\begin{frame}[fragile,shrink]{Descent datum 的完整定义}
设：
$$
F:\mathcal C^{op}\to\mathbf{Cat}
$$

是一个 contravariant pseudofunctor，采用：
$$
c_{f,g}:g^*f^*\Rightarrow(fg)^*
$$

和：
$$
u_X:\operatorname{Id}_{F(X)}
\Rightarrow(\operatorname{id}_X)^*
$$

的 convention。

给定：
$$
S\hookrightarrow y_X
$$

一个 \textbf{\(S\)-descent datum} 是以下数据：

\textbf{(D1) Local objects}

对每一个：
$$
f:Y\to X\in S(Y),
$$

给出：
$$
\boxed{
A_f\in F(Y).
}
$$

\textbf{(D2) Transition isomorphisms}

对每一个：
$$
f:Y\to X\in S(Y),
\qquad
g:Z\to Y,
$$

给出：
$$
\boxed{
\theta_{f,g}:
g^*A_f
\xrightarrow{\cong}
A_{fg}.
}
$$

\textbf{(D3) Identity coherence}

对每一个：
$$
f:Y\to X\in S(Y),
$$

满足：
$$
\boxed{
\theta_{f,\operatorname{id}_Y}
\circ
u_Y(A_f)
=
\operatorname{id}_{A_f}.
}
$$

\textbf{(D4) Cocycle coherence}

对于：
$$
X\xleftarrow{f}Y\xleftarrow{g}Z\xleftarrow{h}W,
$$

满足：
$$
\boxed{
\theta_{fg,h}
\circ
h^*(\theta_{f,g})
=
\theta_{f,gh}
\circ
c_{g,h}(A_f).
}
$$

这四项就是完整的 descent datum。
\end{frame}

\begin{frame}[fragile,shrink]{Global object 如何产生 descent datum？}
这是理解整个定义最重要的一步。

给定：
$$
A\in F(X),
$$

定义：
$$
\boxed{
A_f:=f^*A.
}
$$

对于：
$$
f:Y\to X,
\qquad
g:Z\to Y,
$$

定义：
$$
\boxed{
\theta_{f,g}:=c_{f,g}(A).
}
$$

因为：
$$
c_{f,g}(A):
g^*f^*A
\overset{\cong}{\longrightarrow}
(fg)^*A,
$$

恰好就是：
$$
g^*A_f
\overset{\cong}{\longrightarrow}
A_{fg}.
$$

所以每一个 global object：
$$
A\in F(X)
$$

都会产生一个 descent datum：
$$
\boxed{
\operatorname{Desc}(A).
}
$$

于是得到 canonical functor：
$$
\boxed{
F(X)
\longrightarrow
\operatorname{Desc}_F(S).
}
$$
\end{frame}

\begin{frame}[fragile]{为什么 cocycle condition 正好成立？}
因为 descent datum 的 cocycle condition：
$$
\theta_{fg,h}
\circ
h^*(\theta_{f,g})
=
\theta_{f,gh}
\circ
c_{g,h}(A_f)
$$

在 global descent datum 中变成：
$$
c_{fg,h}(A)
\circ
h^*\bigl(c_{f,g}(A)\bigr)
=
c_{f,gh}(A)
\circ
c_{g,h}(f^*A).
$$

这正是 pseudofunctor 本身的 associativity coherence：
$$
\boxed{
c_{fg,h}\circ(h^**c_{f,g})
=
c_{f,gh}\circ(c_{g,h}*f^*).
}
$$

所以：
$$
\boxed{
\text{descent cocycle}
\quad\text{来自}\quad
\text{pseudofunctor coherence}.
}
$$

这其实是整个结构最漂亮的地方。
\end{frame}

\begin{frame}[fragile,shrink]{Descent category}
最后，如果有两个 descent data：
$$
(A_f,\theta_{f,g})
$$

和：
$$
(B_f,\eta_{f,g}),
$$

它们之间的 morphism 是一族：
$$
\boxed{
\phi_f:A_f\to B_f
}
$$

其中：
$$
\phi_f\in F(Y)
$$

（这里 \(f:Y\to X\)），并且满足：
$$
\boxed{
\eta_{f,g}
\circ
g^*(\phi_f)
=
\phi_{fg}
\circ
\theta_{f,g}.
}
$$

即交换图：
$$
\boxed{
\begin{array}{ccc}
g^*A_f
&
\xrightarrow{g^*(\phi_f)}
&
g^*B_f
\\[2mm]
\downarrow{\theta_{f,g}}
&&
\downarrow{\eta_{f,g}}
\\[2mm]
A_{fg}
&
\xrightarrow{\phi_{fg}}
&
B_{fg}.
\end{array}
}
$$
\end{frame}

\begin{frame}[fragile,shrink]
这些 descent data 和 morphisms 构成 category：
$$
\boxed{
\operatorname{Desc}_F(S).
}
$$

于是整个结构最终可以浓缩为：
$$
\boxed{
F(X)
\longrightarrow
\operatorname{Desc}_F(S),
}
$$

其中：
$$
\boxed{
A
\longmapsto
\left(
A_f=f^*A,\;
\theta_{f,g}=c_{f,g}(A)
\right).
}
$$

而：
$$
\boxed{
\text{\(F\) 是 stack}
\iff
F(X)\longrightarrow\operatorname{Desc}_F(S)
\text{ 对每个 covering sieve \(S\) 都是 equivalence}.
}
$$
\begin{markdown}
这套符号下，**pseudofunctor 的 $c_{f,g},u_X$** 与 **descent datum 的 $\theta_{f,g}$** 的关系就完全统一了：对于一个 global object $A$，它的 descent datum 的 $\theta$ **就是** pseudofunctor 的 composition constraint $c$ 在 $A$ 上的 component。
\end{markdown}
\end{frame}

\begin{frame}
The Grothendieck construction associates to a contravariant pseudofunctor a fibered category, and conversely, each fibered category is induced by some contravariant pseudofunctor. Because of this, a contravariant pseudofunctor, which is a category-valued presheaf, is often also called a prestack (a stack minus effective descent).\cite{wiki_pseudofunctor}
\end{frame}

% 参考文献帧，长列表自动分页
\begin{frame}[allowframebreaks]{参考文献}
\begin{thebibliography}{99}

\bibitem{wiki_pseudofunctor}
\url{https://en.wikipedia.org/wiki/Pseudo-functor}

\end{thebibliography}
\end{frame}
\end{document}